% =====================================================================
\chapter{Series de potencias y series de Taylor}
\label{tema:6}
% =====================================================================

Una serie de potencias es un «polinomio de grado infinito». Dentro de su intervalo de convergencia se comporta como un polinomio: puede derivarse e integrarse término a término. Las funciones elementales $e^{x}$, $\sin x$, $\cos x$, $\ln(1+x)$ y $(1+x)^{a}$ son sumas de series de potencias, sus series de Taylor, y esa representación permite calcularlas con la precisión deseada, obtener límites y definir el número $e$. Este tema estudia el radio de convergencia, las operaciones con series y los desarrollos fundamentales, con demostración de que convergen a la función. Se siguen los capítulos~20 y~24 de Spivak~\cite{spivak}, los capítulos~7 y~11 de Apostol~\cite{apostol}, las §6.1 a~6.4 de~\cite{openstax2}, las §10.4 y~10.5 de Strang~\cite{strang} y las sesiones~97 a~100 del curso 18.01SC~\cite{mit1801}.

\section{Series de potencias y radio de convergencia}

\begin{definicion}[Serie de potencias]\label{def:serie-potencias}
Una \emph{serie de potencias} centrada en $c\in\R$ es una serie de la forma $\sum_{n=0}^{\infty}a_n(x-c)^{n}$, donde $(a_n)$ es una sucesión de números reales y $x$ es una variable real. Su \emph{dominio de convergencia} es el conjunto de los $x$ para los que converge.
\end{definicion}

Toda serie de potencias converge en $x=c$, donde vale $a_0$. El resultado siguiente muestra que el dominio de convergencia es siempre un intervalo centrado en $c$.

\begin{lema}[Lema de Abel]\label{lem:abel}
Si $\sum a_n(x_0-c)^{n}$ converge para algún $x_0\neq c$, entonces $\sum a_n(x-c)^{n}$ converge absolutamente para todo $x$ con $\abs{x-c}<\abs{x_0-c}$.
\end{lema}

\begin{teorema}[Radio de convergencia]\label{teo:radio}
Para cada serie de potencias $\sum a_n(x-c)^{n}$ existe $R\in[0,\infty]$, su \emph{radio de convergencia}, tal que la serie converge absolutamente si $\abs{x-c}<R$ y diverge si $\abs{x-c}>R$. Si existe $\lim_{n\to\infty}\Bigl|\frac{a_n}{a_{n+1}}\Bigr|$ (finito o infinito), ese límite es $R$.
\end{teorema}

El intervalo $(c-R,c+R)$ es el \emph{intervalo de convergencia}. En los extremos $x=c\pm R$ el teorema no dice nada, y cada caso se estudia aparte con los criterios del \cref{tema:5}.

\addtocounter{definicion}{1}%

\section{Derivación e integración término a término}

\begin{teorema}[Derivación e integración de series de potencias]\label{teo:termino-a-termino}
Sea $f(x)=\sum_{n=0}^{\infty}a_n(x-c)^{n}$ con radio de convergencia $R>0$. Las series
\[
\sum_{n=1}^{\infty}n\,a_n(x-c)^{n-1}\qquad\text{y}\qquad\sum_{n=0}^{\infty}\frac{a_n}{n+1}(x-c)^{n+1}
\]
tienen el mismo radio $R$; $f$ es derivable en $(c-R,c+R)$, con $f'$ igual a la primera serie, y $\int_c^{x}f(t)\dif t$ es igual a la segunda para todo $x$ de ese intervalo.
\end{teorema}

La demostración usa la convergencia uniforme de la serie en los intervalos cerrados contenidos en $(c-R,c+R)$ y puede consultarse en~\cite[cap.~24]{spivak} o \cite[cap.~11]{apostol}. Aplicando el teorema repetidamente, $f$ tiene derivadas de todos los órdenes en el intervalo de convergencia.

\begin{corolario}[Unicidad de los coeficientes]\label{cor:coeficientes}
Si $f(x)=\sum_{n=0}^{\infty}a_n(x-c)^{n}$ con $R>0$, entonces $a_n=\frac{f^{(n)}(c)}{n!}$ para todo $n$. En particular, dos series de potencias que coinciden en un entorno de $c$ tienen los mismos coeficientes.
\end{corolario}

\addtocounter{definicion}{1}%

\section{Series de Taylor}

\begin{definicion}[Serie de Taylor]\label{def:serie-taylor}
Sea $f$ con derivadas de todos los órdenes en un intervalo abierto que contiene a $c$. Su \emph{serie de Taylor} en $c$ es $\sum_{n=0}^{\infty}\frac{f^{(n)}(c)}{n!}(x-c)^{n}$. Para $c=0$ se llama también \emph{serie de Maclaurin}. La suma parcial de grado $n$ es el polinomio de Taylor $T_nf(x)$, y $R_nf(x)=f(x)-T_nf(x)$ es el \emph{resto}.
\end{definicion}

Por el \cref{cor:coeficientes}, si $f$ es suma de alguna serie de potencias centrada en $c$, esa serie es su serie de Taylor. Pero la serie de Taylor de $f$ puede converger y no tener por suma $f$ (\cref{eje:6-no-analitica}). La herramienta para decidirlo es el teorema de Taylor del módulo~0.5.

\begin{teorema}[Taylor con resto de Lagrange]\label{teo:taylor-lagrange}
Sea $f$ con $n+1$ derivadas en un intervalo abierto $I$ que contiene a $c$. Para cada $x\in I$ existe $\xi$ entre $c$ y $x$ tal que
\[
R_nf(x)=\frac{f^{(n+1)}(\xi)}{(n+1)!}(x-c)^{n+1}.
\]
\end{teorema}

\begin{proposicion}\label{prop:taylor-converge}
La serie de Taylor de $f$ en $c$ converge a $f(x)$ si y solo si $R_nf(x)\to0$ cuando $n\to\infty$.
\end{proposicion}

\section{Los desarrollos fundamentales}

\begin{teorema}[Exponencial]\label{teo:exp}
Para todo $x\in\R$, $\displaystyle e^{x}=\sum_{n=0}^{\infty}\frac{x^{n}}{n!}=1+x+\frac{x^{2}}{2}+\frac{x^{3}}{6}+\cdots$.
\end{teorema}

Del mismo modo, como las derivadas de $\sin$ y $\cos$ son $\pm\sin$ y $\pm\cos$, acotadas por~1, sus restos cumplen $\abs{R_n(x)}\le\frac{\abs x^{n+1}}{(n+1)!}\to0$, y para todo $x\in\R$ (\cref{eje:6-seno-dem})
\[
\sin x=\sum_{n=0}^{\infty}\frac{(-1)^{n}x^{2n+1}}{(2n+1)!}=x-\frac{x^{3}}{6}+\cdots,\qquad
\cos x=\sum_{n=0}^{\infty}\frac{(-1)^{n}x^{2n}}{(2n)!}=1-\frac{x^{2}}{2}+\cdots.
\]

\begin{teorema}[Logaritmo]\label{teo:ln}
Para todo $x\in(-1,1]$, $\displaystyle\ln(1+x)=\sum_{n=1}^{\infty}\frac{(-1)^{n+1}}{n}x^{n}$. En particular, $\displaystyle\ln2=1-\frac12+\frac13-\frac14+\cdots$.
\end{teorema}

La demostración cubre el extremo $x=1$, que el \cref{teo:termino-a-termino} no alcanza, y da además una cota explícita del error. La serie no sirve para $x>1$, porque allí su radio de convergencia, $R=1$, ya no la hace converger, aunque $\ln(1+x)$ esté definido.

\begin{teorema}[Serie binomial]\label{teo:binomial}
Sea $a\in\R$. Para todo $x\in(-1,1)$,
\[
(1+x)^{a}=\sum_{n=0}^{\infty}\binom{a}{n}x^{n},\qquad\binom{a}{n}=\frac{a(a-1)\cdots(a-n+1)}{n!},\quad\binom{a}{0}=1 .
\]
Si $a\in\N_0$, la serie es finita y se reduce al binomio de Newton.
\end{teorema}

Si $a\notin\N_0$, ningún coeficiente es nulo, $\bigl|\binom{a}{n}\big/\binom{a}{n+1}\bigr|=\frac{n+1}{\abs{a-n}}\to1$ y el radio es~1. La demostración de que la suma es $(1+x)^{a}$ se pide en el \cref{eje:6-binomial-dem}. Por ejemplo, con $a=\frac12$: $\sqrt{1+x}=1+\frac x2-\frac{x^{2}}{8}+\frac{x^{3}}{16}-\cdots$.

\begin{table}[h]
\centering
\small
\renewcommand{\arraystretch}{1.25}
\begin{tabular}{@{}>{$\displaystyle}l<{$}>{$\displaystyle}l<{$}>{$}l<{$}@{}}
\toprule
\text{Función} & \text{Serie de Maclaurin} & \text{Validez} \\
\midrule
\frac{1}{1-x} & \sum_{n\ge0}x^{n} & (-1,1) \\
e^{x} & \sum_{n\ge0}\frac{x^{n}}{n!} & \R \\
\sin x & \sum_{n\ge0}\frac{(-1)^{n}x^{2n+1}}{(2n+1)!} & \R \\
\cos x & \sum_{n\ge0}\frac{(-1)^{n}x^{2n}}{(2n)!} & \R \\
\ln(1+x) & \sum_{n\ge1}\frac{(-1)^{n+1}x^{n}}{n} & (-1,1] \\
\arctan x & \sum_{n\ge0}\frac{(-1)^{n}x^{2n+1}}{2n+1} & [-1,1] \\
(1+x)^{a} & \sum_{n\ge0}\binom{a}{n}x^{n} & (-1,1) \\
\bottomrule
\end{tabular}
\caption{Desarrollos fundamentales. La validez en $x=\pm1$ de la serie del arcotangente se prueba en el \cref{eje:6-arctan}.}
\label{tab:desarrollos}
\end{table}

\section{El número \texorpdfstring{$e$}{e}}

El número $e$ es la base de la exponencial, $e=e^{1}$. Por el \cref{teo:exp},
\[
e=\sum_{n=0}^{\infty}\frac{1}{n!}=1+1+\frac12+\frac16+\frac{1}{24}+\cdots\approx2{,}718281828 .
\]

\begin{proposicion}\label{prop:e-limite}
$\displaystyle\lim_{n\to\infty}\Bigl(1+\frac1n\Bigr)^{n}=e$.
\end{proposicion}

\begin{teorema}[El número $e$ es irracional]\label{teo:e-irracional}
Sea $s_q=\sum_{k=0}^{q}\frac{1}{k!}$. Para todo $q\in\N$, $0<e-s_q<\frac{1}{q!\,q}$. En consecuencia, $e\notin\Q$.
\end{teorema}

La cota del teorema también es práctica: con $q=10$, $s_{10}$ aproxima $e$ con error menor que $\frac{1}{10!\cdot10}\approx2{,}8\cdot10^{-8}$.

\addtocounter{definicion}{1}%

\addtocounter{definicion}{1}%

